\chapter{Lên Thành Phố, Sống Xa Nhà, Tự Lo Mọi Thứ}
\markboth{Lên Thành Phố, Sống Xa Nhà, Tự Lo Mọi Thứ}{Moving Away and Taking Care of Everything}

\section{Sống Xa Nhà Là Một Cú Lớn Lên Rất Thật}
\begin{truyen}{Sống Xa Nhà Là Một Cú Lớn Lên Rất Thật}{Living Away from Home Is a Very Real Shock of Growth}
\chuhoa{K}{hi còn ở nhà, nhiều việc nhỏ được ai đó âm thầm gánh hộ em.}
\textit{(When you live at home, many small burdens are quietly carried for you.)}

Ra ở riêng rồi, em mới thấy bữa ăn, tiền điện, ga gối, bệnh vặt, và cả sự im lặng cũng đều có trọng lượng.
\textit{(Once you live alone, meals, utility bills, laundry, small illness, and even silence all gain weight.)}

Sống xa nhà không chỉ là đổi địa chỉ.
\textit{(Living away from home is not only changing address.)}

Nó là đổi vai trò.
\textit{(It is changing roles.)}

Từ người được chăm, em dần thành người tự chăm mình.
\textit{(From being cared for, you slowly become the one who cares for yourself.)}

\end{truyen}

\begin{baihoc}
Xa nhà làm người ta mệt hơn, nhưng cũng thật hơn. Nó ép em chạm vào phần đời sống mà trước đó em chỉ đi qua.
\textit{(Living away from home is harder, but also more real. It forces you to touch parts of life you once only passed by.)}
\end{baihoc}

\begin{ghinhoanh}
\textbf{Short English to remember:}

Living away from home is harder, but also more real.

It forces you to touch parts of life you once only passed by.

\end{ghinhoanh}

\begin{tuhocnhanh}
\textbf{Cau hoi tu hoc / Self-study questions}

1. Van de chinh la gi? \textit{What is the main problem?}

2. Bai hoc thuc te la gi? \textit{What is the practical lesson?}

3. Mot cau tieng Anh em muon nho la cau nao? \textit{Which English sentence do you want to remember?}
\end{tuhocnhanh}
\ngancach

\section{Bệnh Một Mình Mới Hiểu Nhớ Nhà Là Gì}
\begin{truyen}{Bệnh Một Mình Mới Hiểu Nhớ Nhà Là Gì}{You Understand Homesickness When You Get Sick Alone}
\chuhoa{C}{ó những ngày em sốt, nằm trong phòng trọ, và không có ai hỏi đã ăn chưa.}
\textit{(There are days when you have a fever in a rented room and no one asks whether you have eaten.)}

Lúc đó nỗi nhớ nhà không còn là cảm xúc mơ hồ.
\textit{(Then homesickness is no longer vague.)}

Nó thành một cảm giác rất cụ thể trong cổ họng.
\textit{(It becomes something very physical in your throat.)}

Người trẻ hay nghĩ mình mạnh cho đến khi phải tự lo lúc cơ thể yếu nhất.
\textit{(Young people often think they are strong until they must care for themselves at their weakest.)}

Những ngày như vậy dạy em giá trị của gia đình theo cách không cuốn sách nào dạy hết được.
\textit{(Days like that teach the value of family in a way no book can fully teach.)}

\end{truyen}

\begin{baihoc}
Sống xa nhà khiến em hiểu tình thương không phải điều hiển nhiên. Nó là một thứ rất cụ thể khi em thiếu nó.
\textit{(Living away from home teaches that care is not abstract. You feel exactly what it is when you do not have it.)}
\end{baihoc}

\begin{ghinhoanh}
\textbf{Short English to remember:}

Living away from home teaches that care is not abstract.

You feel exactly what it is when you do not have it.

\end{ghinhoanh}

\begin{tuhocnhanh}
\textbf{Cau hoi tu hoc / Self-study questions}

1. Van de chinh la gi? \textit{What is the main problem?}

2. Bai hoc thuc te la gi? \textit{What is the practical lesson?}

3. Mot cau tieng Anh em muon nho la cau nao? \textit{Which English sentence do you want to remember?}
\end{tuhocnhanh}
\ngancach

\section{Thành Phố Không Ác, Nhưng Không Rảnh Dịu Dàng Với Em}
\begin{truyen}{Thành Phố Không Ác, Nhưng Không Rảnh Dịu Dàng Với Em}{The City Is Not Cruel, But It Is Not Gentle for Free}
\chuhoa{T}{hành phố lớn không ghét em.}
\textit{(A big city does not hate you.)}

Nó chỉ rất bận.
\textit{(It is simply very busy.)}

Nếu em chậm, thiếu tỉnh, hay thiếu tiền, thành phố sẽ cho em thấy giá của những điều đó rất nhanh.
\textit{(If you are slow, careless, or short of money, the city shows the price of those things very quickly.)}

Điều đó làm nhiều người trẻ tưởng mình bị cuộc đời nhắm vào.
\textit{(That makes many young people feel personally attacked by life.)}

Thực ra không phải.
\textit{(In truth, no.)}

Đô thị chỉ vận hành nhanh và lạnh hơn làng quê hay mái nhà cũ của em.
\textit{(Urban life simply moves faster and colder than your old home.)}

\end{truyen}

\begin{baihoc}
Muốn sống được ở thành phố, người trẻ cần tỉnh táo, tự lập, và biết giữ nhịp. Sự ngây thơ ở đây rất đắt.
\textit{(To live well in the city, young people need alertness, self-reliance, and rhythm. Naivety becomes expensive there.)}
\end{baihoc}

\begin{ghinhoanh}
\textbf{Short English to remember:}

To live well in the city, young people need alertness, self-reliance, and rhythm.

Naivety becomes expensive there.

\end{ghinhoanh}

\begin{tuhocnhanh}
\textbf{Cau hoi tu hoc / Self-study questions}

1. Van de chinh la gi? \textit{What is the main problem?}

2. Bai hoc thuc te la gi? \textit{What is the practical lesson?}

3. Mot cau tieng Anh em muon nho la cau nao? \textit{Which English sentence do you want to remember?}
\end{tuhocnhanh}
\ngancach

\section{Một Căn Phòng Nhỏ Cũng Có Thể Dạy Em Rất Nhiều}
\begin{truyen}{Một Căn Phòng Nhỏ Cũng Có Thể Dạy Em Rất Nhiều}{A Small Room Can Teach You a Lot}
\chuhoa{N}{hiều người trẻ bắt đầu đời sống trưởng thành trong một căn phòng rất nhỏ.}
\textit{(Many young adults begin adult life in a very small room.)}

Nhỏ nhưng trong đó có đủ bài học: cách giữ sạch, cách tiết kiệm, cách sống gọn, và cách chịu một mình.
\textit{(It is small, but it holds many lessons: how to keep things clean, save money, live simply, and endure solitude.)}

Không gian chật đôi khi lại ép con người học những điều rất căn bản.
\textit{(A tight space can force a person to learn very basic truths.)}

Có những người sau này có nhà rộng hơn, nhưng trưởng thành nhất lại bắt đầu từ căn phòng bé nhất.
\textit{(Some later live in bigger homes, yet their real growth began in the smallest room.)}

\end{truyen}

\begin{baihoc}
Đừng khinh giai đoạn sống chật. Nhiều nền nếp tốt của người lớn bắt đầu từ những điều rất chật nhưng rất thật.
\textit{(Do not despise the cramped phase of life. Many strong adult habits begin in conditions that are small but very real.)}
\end{baihoc}

\begin{ghinhoanh}
\textbf{Short English to remember:}

Do not despise the cramped phase of life.

Many strong adult habits begin in conditions that are small but very real.

\end{ghinhoanh}

\begin{tuhocnhanh}
\textbf{Cau hoi tu hoc / Self-study questions}

1. Van de chinh la gi? \textit{What is the main problem?}

2. Bai hoc thuc te la gi? \textit{What is the practical lesson?}

3. Mot cau tieng Anh em muon nho la cau nao? \textit{Which English sentence do you want to remember?}
\end{tuhocnhanh}
\ngancach

\section{Gọi Về Nhà Không Phải Là Yếu}
\begin{truyen}{Gọi Về Nhà Không Phải Là Yếu}{Calling Home Is Not Weakness}
\chuhoa{C}{ó người trẻ sống xa nhà và cố tỏ ra cứng đến mức ngay cả khi mệt cũng không muốn gọi về.}
\textit{(Some young people live far from home and act so tough that they refuse to call even when exhausted.)}

Họ sợ mình sẽ yếu đi nếu thừa nhận rằng mình nhớ nhà hoặc cần một lời hỏi han.
\textit{(They fear becoming weak if they admit that they miss home or need a kind voice.)}

Nhưng trưởng thành không có nghĩa là cắt hết nhu cầu được yêu thương.
\textit{(But growing up does not mean cutting away the need to be loved.)}

Biết tự lập và vẫn giữ được sợi dây với gia đình là một kiểu trưởng thành lành mạnh hơn nhiều.
\textit{(Being self-reliant while still keeping a bond with family is a far healthier form of adulthood.)}

\end{truyen}

\begin{baihoc}
Tự lập không đồng nghĩa với cô lập. Người mạnh không phải người không cần ai, mà là người biết cần ai cho đúng.
\textit{(Self-reliance is not self-isolation. A strong person is not one who needs no one, but one who knows whom to lean on wisely.)}
\end{baihoc}

\begin{ghinhoanh}
\textbf{Short English to remember:}

Self-reliance is not self-isolation.

A strong person is not one who needs no one, but one who knows whom to lean on wisely.

\end{ghinhoanh}

\begin{tuhocnhanh}
\textbf{Cau hoi tu hoc / Self-study questions}

1. Van de chinh la gi? \textit{What is the main problem?}

2. Bai hoc thuc te la gi? \textit{What is the practical lesson?}

3. Mot cau tieng Anh em muon nho la cau nao? \textit{Which English sentence do you want to remember?}
\end{tuhocnhanh}
\ngancach

\section{Tết Đầu Tiên Xa Nhà}
\begin{truyen}{Tết Đầu Tiên Xa Nhà}{The First New Year Away from Home}
\chuhoa{C}{ó những ngày lễ đầu tiên xa nhà khiến căn phòng trọ bỗng rộng và lạnh hơn hẳn.}
\textit{(The first holiday away from home can make a rented room suddenly feel wider and colder.)}

Nhiều người trẻ tưởng mình mạnh cho đến khi những ngày sum vầy đến mà mình không có mặt.
\textit{(Many think they are strong until reunion days come and they are not there.)}

Khi đó nhớ nhà không còn là câu nói mơ hồ.
\textit{(Then homesickness is no longer a vague phrase.)}

Nó là cảm giác rất rõ trong ngực.
\textit{(It becomes a clear feeling in the chest.)}

Những ngày như vậy dạy em biết quý mái nhà cũ theo cách trưởng thành hơn.
\textit{(Days like that teach you to value home in a more adult way.)}

\end{truyen}

\begin{baihoc}
Xa nhà không chỉ dạy tự lập. Nó còn dạy biết ơn.
\textit{(Living away teaches not only independence, but also gratitude.)}
\end{baihoc}

\begin{ghinhoanh}
\textbf{Short English to remember:}

Living away teaches not only independence, but also gratitude.

\end{ghinhoanh}

\begin{tuhocnhanh}
\textbf{Cau hoi tu hoc / Self-study questions}

1. Van de chinh la gi? \textit{What is the main problem?}

2. Bai hoc thuc te la gi? \textit{What is the practical lesson?}

3. Mot cau tieng Anh em muon nho la cau nao? \textit{Which English sentence do you want to remember?}
\end{tuhocnhanh}
\ngancach

\section{Tự Nấu Cho Mình}
\begin{truyen}{Tự Nấu Cho Mình}{Cooking for Yourself}
\chuhoa{L}{ần đầu tự nấu một bữa tử tế cho mình, nhiều người mới hiểu sự chăm sóc không tự từ trên trời rơi xuống.}
\textit{(The first time someone cooks a proper meal for themselves, they understand care does not fall from the sky.)}

Trước đó họ tưởng bữa cơm ở nhà là điều tự nhiên sẵn có.
\textit{(Before that, they thought family meals were simply there by nature.)}

Đến khi chính tay lo từng món nhỏ, họ mới thấy đằng sau sự bình thường là rất nhiều công sức.
\textit{(When they must handle each small step, they see how much labor sits behind normal comfort.)}

Một căn bếp nhỏ đôi khi dạy em về tình thương mạnh hơn nhiều bài giảng.
\textit{(A small kitchen can teach you about love more deeply than many lectures.)}

\end{truyen}

\begin{baihoc}
Tự lo bữa ăn làm người trẻ hiểu chăm sóc là hành động cụ thể.
\textit{(Preparing your own meal teaches that care is a concrete action.)}
\end{baihoc}

\begin{ghinhoanh}
\textbf{Short English to remember:}

Preparing your own meal teaches that care is a concrete action.

\end{ghinhoanh}

\begin{tuhocnhanh}
\textbf{Cau hoi tu hoc / Self-study questions}

1. Van de chinh la gi? \textit{What is the main problem?}

2. Bai hoc thuc te la gi? \textit{What is the practical lesson?}

3. Mot cau tieng Anh em muon nho la cau nao? \textit{Which English sentence do you want to remember?}
\end{tuhocnhanh}
\ngancach

\section{Căn Phòng Và Sự Cô Đơn}
\begin{truyen}{Căn Phòng Và Sự Cô Đơn}{A Room and Its Loneliness}
\chuhoa{C}{ó những đêm phòng trọ rất yên nhưng lòng người lại rất ồn.}
\textit{(Some nights the rented room is very quiet while the mind is very loud.)}

Nhiều người tưởng sống riêng chỉ cần lo vật chất là đủ.
\textit{(Many think living alone is only about handling material needs.)}

Nhưng cô đơn cũng là một hóa đơn tinh thần mà người sống xa nhà phải học cách trả.
\textit{(But loneliness is also a mental bill that people away from home must learn to pay.)}

Biết tự làm đầy đời mình bằng nề nếp tốt và kết nối lành mạnh là một phần của trưởng thành.
\textit{(Knowing how to fill your life with healthy routines and healthy connections is part of maturity.)}

\end{truyen}

\begin{baihoc}
Sống xa nhà không chỉ thử ví tiền. Nó thử cả nội lực bên trong.
\textit{(Living away tests not only your wallet, but also your inner strength.)}
\end{baihoc}

\begin{ghinhoanh}
\textbf{Short English to remember:}

Living away tests not only your wallet, but also your inner strength.

\end{ghinhoanh}

\begin{tuhocnhanh}
\textbf{Cau hoi tu hoc / Self-study questions}

1. Van de chinh la gi? \textit{What is the main problem?}

2. Bai hoc thuc te la gi? \textit{What is the practical lesson?}

3. Mot cau tieng Anh em muon nho la cau nao? \textit{Which English sentence do you want to remember?}
\end{tuhocnhanh}
\ngancach

\section{Nhớ Đường Hỏi Giúp}
\begin{truyen}{Nhớ Đường Hỏi Giúp}{Knowing How to Ask for Help}
\chuhoa{C}{ó người lên thành phố mang theo sĩ diện nên có khó cũng không dám hỏi ai.}
\textit{(Some move to the city carrying so much pride that they do not ask for help even in real difficulty.)}

Họ tưởng tự lập nghĩa là phải tự làm hết mọi thứ một mình.
\textit{(They think independence means doing everything alone.)}

Nhưng người trưởng thành khỏe mạnh biết lúc nào cần tự lo và lúc nào cần mở lời đúng người.
\textit{(But healthy adults know when to handle things alone and when to ask the right person for help.)}

Xin giúp đỡ không làm em nhỏ đi.
\textit{(Asking for help does not make you smaller.)}

Nó chỉ giúp em bớt ngu ngơ trước giới hạn của mình.
\textit{(It makes you less foolish about your limits.)}

\end{truyen}

\begin{baihoc}
Tự lập không phải cô lập. Đó là biết dựa đúng chỗ đúng lúc.
\textit{(Self-reliance is not isolation. It is knowing where to lean at the right moment.)}
\end{baihoc}

\begin{ghinhoanh}
\textbf{Short English to remember:}

Self-reliance is not isolation.

It is knowing where to lean at the right moment.

\end{ghinhoanh}

\begin{tuhocnhanh}
\textbf{Cau hoi tu hoc / Self-study questions}

1. Van de chinh la gi? \textit{What is the main problem?}

2. Bai hoc thuc te la gi? \textit{What is the practical lesson?}

3. Mot cau tieng Anh em muon nho la cau nao? \textit{Which English sentence do you want to remember?}
\end{tuhocnhanh}
\ngancach

\section{Biết Ơn Cha Mẹ Hơn}
\begin{truyen}{Biết Ơn Cha Mẹ Hơn}{Understanding Parents More Deeply}
\chuhoa{S}{au vài năm sống xa, có người bất chợt thấy mình dịu hơn khi nghĩ về bố mẹ.}
\textit{(After some years away, some people suddenly feel softer when thinking about their parents.)}

Trước đó họ tưởng nhiều điều cha mẹ làm chỉ là nghiêm và khó tính.
\textit{(Before that, they thought many parental choices were only strict or difficult.)}

Nhưng khi tự gánh đời sống, họ mới hiểu nhiều nỗi lo của người lớn không dễ gọi tên.
\textit{(But once they carry life themselves, they understand that many adult worries are hard to describe.)}

Khoảng cách đôi khi làm tình thân rõ hơn chứ không hề mờ đi.
\textit{(Distance sometimes clarifies family love instead of fading it.)}

\end{truyen}

\begin{baihoc}
Lớn lên đẹp là lúc em bớt trách quá khứ và hiểu hơn phần mệt của cha mẹ.
\textit{(A beautiful kind of growing up is blaming the past less and understanding your parents more.)}
\end{baihoc}

\begin{ghinhoanh}
\textbf{Short English to remember:}

A beautiful kind of growing up is blaming the past less and understanding your parents more.

\end{ghinhoanh}

\begin{tuhocnhanh}
\textbf{Cau hoi tu hoc / Self-study questions}

1. Van de chinh la gi? \textit{What is the main problem?}

2. Bai hoc thuc te la gi? \textit{What is the practical lesson?}

3. Mot cau tieng Anh em muon nho la cau nao? \textit{Which English sentence do you want to remember?}
\end{tuhocnhanh}
